%% sockadv.tex
%%
%% Copyright (C) 2004-2018 Simone Piccardi.  Permission is granted to
%% copy, distribute and/or modify this document under the terms of the GNU Free
%% Documentation License, Version 1.1 or any later version published by the
%% Free Software Foundation; with the Invariant Sections being "Un preambolo",
%% with no Front-Cover Texts, and with no Back-Cover Texts.  A copy of the
%% license is included in the section entitled "GNU Free Documentation
%% License".
%%

\chapter{Socket avanzati}
\label{cha:advanced_socket}


Esamineremo in questo capitolo le funzionalità più evolute della gestione dei
socket, le funzioni avanzate, la gestione dei dati urgenti e
\textit{out-of-band} e dei messaggi ancillari, come l'uso come l'uso del I/O
multiplexing (vedi sez.~\ref{sec:file_multiplexing}) con i socket.


\section{Le funzioni di I/O avanzate}
\label{sec:sock_advanced_IO}

Tratteremo in questa sezione le funzioni di I/O più avanzate che permettono di
controllare le funzionalità specifiche della comunicazione dei dati che sono
disponibili con i vari tipi di socket.

\subsection{La funzioni \func{send} e \func{recv}}
\label{sec:net_send_recv}

Da fare


% TODO: note su MSG_ZEROCOPY/SOCK_ZEROCOPY, aggiunte con il kernel 4.14 (e per
% la ricezione con il kernel 4.18, vedi https://lwn.net/Articles/726917/ e
% https://lwn.net/Articles/752300/ 

\subsection{La funzioni \func{sendmsg} e \func{recvmsg}}
\label{sec:net_sendmsg}

Finora abbiamo trattato delle funzioni che permettono di inviare dati sul
socket in forma semplificata. Se infatti si devono semplicemente ...

% TODO trattare anche recvmmsg, introdotta con il kernel 2.3.33, vedi
% http://kernelnewbies.org/Linux_2_6_33 

% TODO trattare anche sendmmsg, introdotta con il kernel 3.0, vedi
% 


\subsection{I messaggi ancillari}
\label{sec:net_ancillary_data}

Quanto è stata attivata l'opzione \const{IP\_RECVERR} il kernel attiva per il
socket una speciale coda su cui vengono inviati tutti gli errori riscontrati.
Questi possono essere riletti usando il flag \const{MSG\_ERRQUEUE}, nel qual
caso sarà passato come messaggio ancillare una struttura di tipo
\struct{sock\_extended\_err} illustrata in
fig.~\ref{fig:sock_extended_err_struct}.

\begin{figure}[!htb]
  \footnotesize \centering
  \begin{minipage}[c]{\textwidth}
    \includestruct{listati/sock_extended_err.h}
  \end{minipage}
  \caption{La struttura \structd{sock\_extended\_err} usata dall'opzione
    \const{IP\_RECVERR} per ottenere le informazioni relative agli errori su
    un socket.}
  \label{fig:sock_extended_err_struct}
\end{figure}

% TODO vedi man cmsg

\subsection{I \textsl{dati urgenti} o \textit{out-of-band}}
\label{sec:TCP_urgent_data}

\itindbeg{out-of-band} 

Una caratteristica particolare dei socket TCP è quella che consente di inviare
all'altro capo della comunicazione una sorta di messaggio privilegiato, che si
richiede che sia trattato il prima possibile. Si fa riferimento a questa
funzionalità come all'invio dei cosiddetti \textsl{dati urgenti} (o
\textit{urgent data}); talvolta essi chiamati anche dati \textit{out-of-band}
poiché, come vedremo più avanti, possono essere letti anche al di fuori del
flusso di dati normale.

Come già accennato in sez.~\ref{sec:file_multiplexing} la presenza di dati
urgenti viene rilevata in maniera specifica sia di \func{select} (con il
\textit{file descriptor set} \param{exceptfds}) che da \func{poll} (con la
condizione \const{POLLRDBAND}).


Le modalità di lettura dei dati urgenti sono due, la prima e più comune
prevede l'uso di \func{recvmsg} con 


% TODO aggiungere pezzo di codice per inviare dati urgenti all'echo server

La seconda modalità di lettura prevede invece l'uso dell'opzione dei socket
\const{SO\_OOBINLINE} (vedi sez.~\ref{sec:sock_generic_options}) che consente
di ricevere i dati urgenti direttamente nel flusso dei dati del socket; in tal
caso però si pone il problema di come distinguere i dati normali da quelli
urgenti. Come già accennato in sez.~\ref{sec:sock_ioctl_IP} a questo scopo si
può usare \func{ioctl} con l'operazione \const{SIOCATMARK}, che consente di
sapere se si è arrivati o meno all'\textit{urgent mark}. 

La procedura allora prevede che, una volta che si sia rilevata la presenza di
dati urgenti, si ripeta la lettura ordinaria dal socket fintanto che
\const{SIOCATMARK} non restituisce un valore diverso da zero; la successiva
lettura restituirà i dati urgenti.


\itindend{out-of-band} 


% TODO trattare sendmmsg aggiunta con il kernel 3.0.
% e gli ICMP socket, vedi  http://lwn.net/Articles/420799/


\section{L'uso dell'I/O non bloccante}
\label{sec:sock_noblok_IO}

Tratteremo in questa sezione le modalità avanzate che permettono di utilizzare
i socket con una comunicazione non bloccante, in modo da 





\subsection{La gestione delle opzioni IP}  
\label{sec:sock_IP_options}  

Abbiamo visto in sez.~\ref{sec:sock_ipv4_options} come di possa usare  
\func{setsockopt} con l'opzione \const{IP\_OPTIONS} per impostare le opzioni  
IP associate per i pacchetti associati ad un socket.  Vedremo qui il
significato di tali opzioni e le modalità con cui esse possono essere
utilizzate ed impostate.

% LocalWords:  socket of multiplexing sez sendmsg recvmsg RECVERR kernel MSG
% LocalWords:  ERRQUEUE sock err fig TCP dell'I setsockopt OPTIONS urgent poll
% LocalWords:  select descriptor exceptfds POLLRDBAND OOBINLINE ioctl all' mark
% LocalWords:  SIOCATMARK

%%% Local Variables: 
%%% mode: latex
%%% TeX-master: "gapil"
%%% End: 
